\documentclass[10pt,a4paper]{amsart}
\usepackage[margin=19mm]{geometry}
\usepackage{amsmath,amssymb,mathtools}
\usepackage[scr=rsfs,cal=euler]{mathalfa}
\usepackage{xfrac}
\usepackage[unicode,hidelinks,bookmarks=false]{hyperref}
\newcommand{\ie}{i.e.\ }
\newcommand{\cf}{cf.\ }
\newcommand{\resp}{respectively\ }
\newcommand{\Subsection}{\subsection}
\providecommand{\nobreakdash}{\nobreak\hbox{-}}
\setlength{\emergencystretch}{2em}
\multlinegap0pt
\usepackage{amssymb}
\usepackage{mathtools}
\usepackage[shortlabels]{enumitem}
\newtheorem{lemma}{Lemma}
\newtheorem{proposition}{Proposition}
\newcommand{\Id}{\mathrm{Id}}
\newcommand{\av}[1]{\langle #1\rangle}
\newcommand{\cD}{\mathcal D}
\newcommand{\eps}{\varepsilon}
\newcommand{\ii}{\mathrm i}
\newcommand{\norm}[1]{\lVert #1\rVert}
\title{Lagrange Stability for Reversible Duffing Equations\\ with Quasi-Periodic Coefficients}
\author{Huining Xue}
\date{Source: arXiv:2607.17068v1 (CC BY 4.0)}
\begin{document}
\maketitle

\noindent\emph{Source and licence.} Lagrange Stability for Reversible Duffing Equations\\ with Quasi-Periodic Coefficients. Version arXiv:2607.17068v1 (CC BY 4.0), distributed by arXiv under the Creative Commons Attribution 4.0 International licence, http://creativecommons.org/licenses/by/4.0/, as recorded in the arXiv licence metadata for this exact version and shown on the abstract page https://arxiv.org/abs/2607.17068v1. Full licence notice: \texttt{licenses/arxiv-2607.17068-CC-BY-4.0.txt}.\par\smallskip

\noindent\textit{Editorial scope.} Counted unit: the article's proposition prop:nf (source0.tex lines 626--664). The article's complete original proof of it is delivered with it (source0.tex lines 666--956, one \texttt{proof} environment of 291 lines). No mathematics is written, repaired or re-derived: the statement and the proof are the article's own lines, in the article's own notation, and no retained line is substituted or re-flowed.  Retained uncounted as the source-local context the proof needs: lines 237--241; lines 277--312; lines 417--420; lines 432--434; lines 460--466; lines 529--536; lines 588--592.  Nothing was omitted from any retained span, so the package carries no omission record.  No retained line contains a citation, so the wrapper carries no bibliography and no bibliography entry.  No retained line contains a figure environment, an image-inclusion command or drawing code, so no image is delivered and the package declares no asset dependency.  Editorial formatting only, in addition: an AMS \texttt{amsart} preamble that restates the article's own declarations and macros used by the retained lines, namely the environments \texttt{lemma}, \texttt{proposition} on separate counters, together with the standard packages those lines need.  The retained environments print in this document as Lemma 1, Proposition 1 in order of appearance, whereas the article prints the same items with the numbering of its own full document.  The complete native source of the article as distributed in the arXiv e-print for this exact version is bundled unchanged as \texttt{sources/205528/source0.tex}.
\begin{equation}\label{eq:AAfield}
 \dot\rho=F_1(\rho,\varphi,\theta),\qquad
 \dot\varphi=\nu(\rho)+F_2(\rho,\varphi,\theta),\qquad
 \dot\theta=\omega,
\end{equation}

\begin{lemma}[Equivariant changes preserve reversibility]
\label{lem:revchange}
Let $z=(r,\psi)$, $S=\operatorname{diag}(1,-1)$, and
\[
 \mathscr G(z,\theta)=(Sz,-\theta).
\]
Suppose the extended vector field
$\widetilde X=(X(z,\theta),\omega)$ satisfies
\begin{equation}\label{eq:revabstract}
 \widetilde X(\mathscr G\xi)
 =-D\mathscr G\,\widetilde X(\xi).
\end{equation}
Let
\[
 \widetilde\Phi(z,\theta)=(\Phi(z,\theta),\theta)
\]
be an analytic diffeomorphism satisfying
\begin{equation}\label{eq:equivariance}
 \widetilde\Phi\circ\mathscr G
 =\mathscr G\circ\widetilde\Phi .
\end{equation}
Then the pull-back vector field
\begin{equation}\label{eq:pullback}
 \widetilde Y(\xi)
 =D\widetilde\Phi(\xi)^{-1}
 \widetilde X(\widetilde\Phi(\xi))
\end{equation}
is reversible with respect to $\mathscr G$.

If $\Phi(z,\theta)=z+W(z,\theta)$ and $W=(U,V)$, condition
\eqref{eq:equivariance} is equivalent to
\begin{equation}\label{eq:UVequiv}
 U(r,-\psi,-\theta)=U(r,\psi,\theta),\qquad
 V(r,-\psi,-\theta)=-V(r,\psi,\theta).
\end{equation}
\end{lemma}

\begin{equation}\label{eq:Fouriertail}
 |\Pi_K^\perp h|_{r,s-\sigma}
 \leq C_d K^d e^{-K\sigma}|h|_{r,s}.
\end{equation}

\begin{equation}\label{eq:K}
 K=c_*\log A.
\end{equation}

\begin{equation}\label{eq:NR}
 |D_{k\ell}(\rho)|\geq
 \begin{cases}
 \gamma(1+|\ell|)^{-d-2},& k=0,\\[2mm]
 \delta A^q(1+|\ell|)^{-d-2},& k\neq0.
 \end{cases}
\end{equation}

\begin{equation}\label{eq:initialNF}
 \begin{aligned}
 \dot\rho&=a_0(\rho,\theta)+F_0(\rho,\varphi,\theta),\\
 \dot\varphi&=\nu(\rho)+b_0(\rho,\theta)
                   +G_0(\rho,\varphi,\theta),\\
 \dot\theta&=\omega,
 \end{aligned}
\end{equation}

\begin{equation}\label{eq:epsilon0}
 \av{F_0}_{\varphi}=\av{G_0}_{\varphi}=0,\qquad
 \norm{(F_0,G_0)}_{r_0,s_0/4}\leq
 \eps_0=A^{-\sigma_0}.
\end{equation}

\begin{proposition}[Finite-step normal form]\label{prop:nf}
Let $L>0$.  If $c_*$ in \eqref{eq:K} and then $A$ are sufficiently large,
there exist an integer
\begin{equation}\label{eq:Mchoice}
 M=\min\{j\geq1:\sigma_0+\tfrac12j\kappa_0\geq L+4\}
\end{equation}
and reversible analytic transformations
\[
 \Phi_v=\Id+W_v,\qquad 0\leq v<M,
\]
such that
\[
 \Psi^{(M)}=\Psi_0\circ\Phi_0\circ\cdots\circ\Phi_{M-1}
\]
transforms \eqref{eq:AAfield} into
\begin{equation}\label{eq:NFfinal}
 \begin{aligned}
 \dot r&=a_M(r,\theta)+F_M(r,\psi,\theta),\\
 \dot\psi&=\nu(r)+b_M(r,\theta)+G_M(r,\psi,\theta),\\
 \dot\theta&=\omega,
 \end{aligned}
\end{equation}
where
\begin{equation}\label{eq:finalsmall}
 \norm{(F_M,G_M)}_{r_M,s_M}\leq A^{-L},
\end{equation}
\begin{equation}\label{eq:finalparity}
 \begin{aligned}
 a_M(r,-\theta)&=-a_M(r,\theta),&
 b_M(r,-\theta)&= b_M(r,\theta),\\
 F_M(r,-\psi,-\theta)&=-F_M(r,\psi,\theta),&
 G_M(r,-\psi,-\theta)&= G_M(r,\psi,\theta).
 \end{aligned}
\end{equation}
The final widths satisfy
\[
 r_M\geq\frac12r_0,\qquad s_M\geq\frac18s_0.
\]
\end{proposition}

\begin{proof}
We give the induction in five steps.

\smallskip
\noindent\emph{Step 1: induction hypotheses.}
Set
\begin{equation}\label{eq:widths}
 \sigma_v=\frac{s_0}{2^{v+6}},\qquad
 s_{v+1}=s_v-4\sigma_v,\qquad
 r_{v+1}=r_v-\frac{r_0}{2^{v+3}},
\end{equation}
with $s_0$ in this display replaced initially by the width $s_0/4$
of \eqref{eq:initialNF}.  Suppose that on
$\cD(I_A,r_v,s_v)$ the $v$-th system is
\begin{equation}\label{eq:vthsystem}
 \dot z=N_v(z,\theta)+R_v(z,\theta),\qquad
 \dot\theta=\omega,
\end{equation}
where $z=(r,\psi)$,
\begin{equation}\label{eq:NvRv}
 N_v=
 \binom{a_v(r,\theta)}
       {\nu(r)+b_v(r,\theta)},\qquad
 R_v=\binom{F_v(r,\psi,\theta)}
           {G_v(r,\psi,\theta)},
\end{equation}
\[
 \av{R_v}_{\psi}=0,\qquad
 \norm{R_v}_{r_v,s_v}\leq\eps_v,
\]
and, with $\overline N_v=(a_v,b_v)^T$,
\begin{equation}\label{eq:averageTameV}
 \Gamma_A\left(
 |D_z\overline N_v|+
 r_v|D_z^2\overline N_v|
 \right)\leq2A^{-\kappa_0}.
\end{equation}
and the parity relations in \eqref{eq:finalparity} hold with $M$ replaced by
$v$.  For $v=0$, these assertions are \eqref{eq:initialNF}--\eqref{eq:epsilon0}.

\smallskip
\noindent\emph{Step 2: the homological equation.}
Let
\[
 R_v^{\leq}=\Pi_KR_v,\qquad R_v^>=\Pi_K^\perp R_v.
\]
We solve the vector homological equation
\begin{equation}\label{eq:homological}
 \mathcal H_0W_v=R_v^{\leq},\qquad
 \mathcal H_0W
 =\omega\cdot\partial_\theta W+
 D_zW\,N_0-D_zN_0\,W,\qquad
 N_0=\binom0{\nu(r)},
\qquad \av{W_v}_{\psi}=0.
\end{equation}
Writing $W_v=(U_v,V_v)$ and
$R_v^{\leq}=(F_v^{\leq},G_v^{\leq})$, its solution is
\begin{align}
 U_v(r,\psi,\theta)
 &=\sum_{\substack{|k|+|\ell|\leq K\\k\neq0}}
 \frac{\widehat F_v(r,k,\ell)}
 {\ii D_{k\ell}(r)}
 e^{\ii(k\psi+\langle\ell,\theta\rangle)},\label{eq:Uv}\\
 V_v(r,\psi,\theta)
 &=\sum_{\substack{|k|+|\ell|\leq K\\k\neq0}}
 \left\{
 \frac{\widehat G_v(r,k,\ell)}
 {\ii D_{k\ell}(r)}
 +\frac{\nu'(r)\widehat F_v(r,k,\ell)}
 {(\ii D_{k\ell}(r))^2}
 \right\}
 e^{\ii(k\psi+\langle\ell,\theta\rangle)}.
\label{eq:Vv}
\end{align}
There are no $k=0$ terms because $\av{R_v}_{\psi}=0$.
By \eqref{eq:NR},
\begin{equation}\label{eq:Gamma}
 \Gamma_A=
 C\delta^{-1}A^{-q}K^{d+2}
 \leq CA^{-q}K^{5d+8}.
\end{equation}
Using Cauchy's estimate on the loss $\sigma_v$,
\begin{equation}\label{eq:Westimate}
 \norm{W_v}_{r_v-\frac{r_0}{2^{v+4}},s_v-\sigma_v}
 \leq
 C\Gamma_A\sigma_v^{-(d+2)}\eps_v
 =:\eta_v.
\end{equation}
Moreover, differentiating $D_{k\ell}^{-1}$ gives
\[
 \partial_r D_{k\ell}^{-1}
 =-\frac{k\nu'(r)}{D_{k\ell}(r)^2},
\]
which is included in the weighted norm in \eqref{eq:Westimate}.

The induction hypothesis gives, for real $r$,
\begin{equation}\label{eq:FourierParity}
 \widehat F_v(r,-k,-\ell)=-\widehat F_v(r,k,\ell),
 \qquad
 \widehat G_v(r,-k,-\ell)= \widehat G_v(r,k,\ell).
\end{equation}
Since
\begin{equation}\label{eq:divisorParity}
 D_{-k,-\ell}(r)=-D_{k\ell}(r),
\end{equation}
formulas \eqref{eq:Uv}--\eqref{eq:Vv} imply
\[
 \widehat U_v(r,-k,-\ell)=\widehat U_v(r,k,\ell),
 \qquad
 \widehat V_v(r,-k,-\ell)=-\widehat V_v(r,k,\ell).
\]
Equivalently,
\begin{equation}\label{eq:Wparity}
 U_v(r,-\psi,-\theta)=U_v(r,\psi,\theta),\qquad
 V_v(r,-\psi,-\theta)=-V_v(r,\psi,\theta).
\end{equation}
Therefore $\Phi_v=\Id+W_v$ commutes with
$(r,\psi,\theta)\mapsto(r,-\psi,-\theta)$.  Lemma
\ref{lem:revchange} now shows directly that the vector field after the
$v$-th change is reversible.

\smallskip
\noindent\emph{Step 3: exact transformed vector field.}
Use old coordinates $z=\Phi_v(z_+,\theta)=z_++W_v(z_+,\theta)$.
Then
\[
 \dot z=(I+D_zW_v)\dot z_+
       +\omega\cdot\partial_\theta W_v.
\]
Consequently,
\begin{equation}\label{eq:exactpush}
 X_{v+1}
 =(I+D_zW_v)^{-1}
 \left[
 N_v\circ\Phi_v+R_v\circ\Phi_v
 -\omega\cdot\partial_\theta W_v
 \right].
\end{equation}
Write $N_v=N_0+\overline N_v$.  Subtract
$(I+D_zW_v)N_v$ in \eqref{eq:exactpush} and use
\eqref{eq:homological}.  We obtain
\begin{equation}\label{eq:pushsplit}
 X_{v+1}
 =N_v+\mathcal E_v,
\end{equation}
where
\begin{equation}\label{eq:Ev}
 \mathcal E_v
 =(I+D_zW_v)^{-1}
 \left[
 R_v^>+\mathcal Q_v
 \right].
\end{equation}
Here
\begin{align}
 \mathcal Q_v={}&
 N_0\circ\Phi_v-N_0-D_zN_0W_v
 +\overline N_v\circ\Phi_v-\overline N_v
\notag\\
&-D_zW_v\,\overline N_v
 +R_v\circ\Phi_v-R_v .
\label{eq:Qv}
\end{align}
The Taylor remainders have the exact integral form
\begin{align}
 N_0\circ\Phi_v-N_0-D_zN_0W_v
 &=\int_0^1(1-t)
 D_z^2N_0(z+tW_v)[W_v,W_v]\,dt,
 \label{eq:TaylorN}\\
 R_v\circ\Phi_v-R_v
 &=\int_0^1
 D_zR_v(z+tW_v)W_v\,dt.
\label{eq:TaylorR}
\end{align}
Also,
\begin{equation}\label{eq:TaylorAverage}
 \overline N_v\circ\Phi_v-\overline N_v
 =\int_0^1D_z\overline N_v(z+tW_v)W_v\,dt.
\end{equation}

Define the new average and oscillatory parts by
\begin{equation}\label{eq:newaverage}
 N_{v+1}=N_v+\av{\mathcal E_v}_{\psi},\qquad
 R_{v+1}=\mathcal E_v-\av{\mathcal E_v}_{\psi}.
\end{equation}
Thus $\av{R_{v+1}}_{\psi}=0$ exactly.  More explicitly, reversibility gives
\[
 \mathcal E_{v,1}(r,-\psi,-\theta)
 =-\mathcal E_{v,1}(r,\psi,\theta),\qquad
 \mathcal E_{v,2}(r,-\psi,-\theta)
 = \mathcal E_{v,2}(r,\psi,\theta).
\]
Taking the $\psi$-average yields
\[
 \av{\mathcal E_{v,1}}_\psi(r,-\theta)
 =-\av{\mathcal E_{v,1}}_\psi(r,\theta),\qquad
 \av{\mathcal E_{v,2}}_\psi(r,-\theta)
 = \av{\mathcal E_{v,2}}_\psi(r,\theta).
\]
Consequently both $N_{v+1}$ and $R_{v+1}$ have the required parity, so
\eqref{eq:finalparity} propagates from $v$ to $v+1$.
Furthermore, \eqref{eq:newaverage}, Cauchy's estimate, and
\eqref{eq:recurrence0} below imply
\[
 \Gamma_A\left(
 |D_z\overline N_{v+1}|+
 r_{v+1}|D_z^2\overline N_{v+1}|
 \right)
 \leq
 \Gamma_A\left(
 |D_z\overline N_v|+
 r_v|D_z^2\overline N_v|
 \right)
 +C\Gamma_A\sigma_v^{-2}\eps_v .
\]
For large $A$, \eqref{eq:epsinduction} makes the last term at most
$2^{-v-2}A^{-\kappa_0}$; hence \eqref{eq:averageTameV} also propagates.

\smallskip
\noindent\emph{Step 4: the recurrence.}
Assume $\eta_v\leq c\min\{r_v-r_{v+1},\sigma_v\}$.
Then
\[
 \norm{(I+D_zW_v)^{-1}}\leq2.
\]
Equations \eqref{eq:Fouriertail},
\eqref{eq:Westimate}, \eqref{eq:Ev},
\eqref{eq:TaylorN}, \eqref{eq:TaylorR}, and
\eqref{eq:TaylorAverage} give
\begin{equation}\label{eq:recurrence0}
 \eps_{v+1}
 \leq
 C A^{-\kappa_0}\sigma_v^{-(d+3)}\eps_v
 +C\sigma_v^{-(d+3)}\Gamma_A\eps_v^2
 +C K^d e^{-K\sigma_v}\eps_v.
\end{equation}
Put
\begin{equation}\label{eq:Bv}
 B_v=C\sigma_v^{-(d+3)}K^{5d+8}.
\end{equation}
Since $A^{-q}\leq1$, \eqref{eq:recurrence0} implies
\begin{equation}\label{eq:recurrence}
 \eps_{v+1}
 \leq C_vA^{-\kappa_0}\eps_v
 +B_v\eps_v^2
 +CK^dA^{-c_*\sigma_v}\eps_v.
\end{equation}

Choose $c_*$ so that
\begin{equation}\label{eq:cstarchoice}
 c_*\min_{0\leq v<M}\sigma_v\geq L+\sigma_0+M\kappa_0+10.
\end{equation}
For $A$ sufficiently large, the second term in
\eqref{eq:recurrence} is bounded by $A^{-L-5}$.
An induction now gives
\begin{equation}\label{eq:epsinduction}
 \eps_v\leq A^{-\sigma_0-\frac12v\kappa_0+1},
 \qquad 0\leq v\leq M.
\end{equation}
Indeed, if \eqref{eq:epsinduction} holds at $v$, then
\[
 C_vA^{-\kappa_0}\eps_v+B_v\eps_v^2
 \leq
 C_v(\log A)^{C_v}
 A^{-\sigma_0-(v+1)\kappa_0+2}
 \leq
 \frac12A^{-\sigma_0-\frac12(v+1)\kappa_0+1},
\]
and \eqref{eq:cstarchoice} gives the same bound for the Fourier tail.
The estimate \eqref{eq:Westimate} and \eqref{eq:epsinduction} also verify
the smallness condition on $\eta_v$, closing the induction.

\smallskip
\noindent\emph{Step 5: termination.}
By \eqref{eq:Mchoice} and \eqref{eq:epsinduction},
\[
 \eps_M\leq
 A^{-\sigma_0-\frac12M\kappa_0+1}
 \leq A^{-L-3}\leq A^{-L}.
\]
Moreover, summing \eqref{eq:widths},
\[
 r_M\geq r_0-\sum_{v=0}^{\infty}\frac{r_0}{2^{v+3}}
 \geq\frac34r_0,
\qquad
 s_M\geq\frac{s_0}{4}
 -4\sum_{v=0}^{\infty}\frac{s_0}{2^{v+6}}
 \geq\frac{s_0}{8}.
\]
This proves \eqref{eq:finalsmall} and completes the finite induction.
\end{proof}

\end{document}
